% Authored by Codex and Claude
\documentclass[11pt]{article}

\usepackage[T1]{fontenc}
\usepackage[margin=1in]{geometry}
\usepackage{amsmath,amssymb,amsthm}
\usepackage{array}

\newtheorem{theorem}{Theorem}
\theoremstyle{definition}
\newtheorem{definition}{Definition}

\newcommand{\Fin}{\operatorname{Fin}}
\newcommand{\run}{\operatorname{run}}
\newcommand{\flushOps}{\operatorname{flushOps}}
\newcommand{\sort}{\operatorname{sort}}
\newcommand{\cmp}{\operatorname{cmp}}
\newcommand{\Some}[1]{\texttt{some}(#1)}
\newcommand{\None}{\texttt{none}}
\newcommand{\intereq}{\simeq_{\mathrm{inter}}}
\newcommand{\flusheq}{\simeq_{\mathrm{flush}}}

\begin{document}
\begin{center}
  {\LARGE\bfseries Flush Equivalence Determines Interleaved Behavior\par}
  \vspace{0.35em}
  {\LARGE\bfseries for Stateless PIFO Trees\par}
\end{center}
\vspace{1.2em}

\section*{Definitions}

\begin{definition}[PIFO]
An entry with values in a set \(X\) is a triple
\[
  (x,\rho,\tau)\in X\times\mathbb N\times\mathbb N,
\]
where \(\rho\) is its rank and \(\tau\) is its arrival stamp.  A PIFO is a
finite queue of entries.  A pop removes the entry with lexicographically
least key \((\rho,\tau)\): lower ranks come first, and equal ranks are served
in FIFO order.  A push appends its entry to the queue.
\end{definition}

\begin{definition}[PIFO tree and annotated path]
A PIFO-tree topology is either a leaf or an internal node with a finite
ordered list of child topologies, possibly empty.  A state places a PIFO at
every node.  A leaf PIFO contains packets, while an internal PIFO contains
child indices.

An annotated root-to-leaf path specifies, at every internal node, a child
index and the rank of the reference inserted there, and specifies the
packet's rank at the final leaf.  Pushing a packet with arrival stamp
\(\tau\) inserts entries carrying that same stamp at every node on the path.
A path is valid when its shape and all its child indices fit the topology.
On a path mismatch the formal push operation leaves the mismatched subtree
unchanged; this clause is unreachable for the valid schedulers considered
below.

To pop a tree, pop its root PIFO and, at an internal node, recurse into the
selected child.  The operation is atomic: it either reaches a leaf and
returns a packet together with the updated tree, or returns \(\None\) and
leaves the original state unchanged.
\end{definition}

\begin{definition}[Stateless scheduler and observations]
For \(k\in\mathbb N\), let \(A=\Fin(k)\).  A stateless scheduler \(S\) on
\(A\) consists of a PIFO-tree topology and one fixed annotated path
\(S(a)\) for each \(a\in A\).  Every push of type \(a\) uses this same path,
independently of the state and arrival time.  The scheduler is
\emph{valid} if every path \(S(a)\) is valid.

An operation is either \(\operatorname{push}(a)\) or
\(\operatorname{pop}\).  A run starts with every queue empty and stamps
successive pushes \(1,2,\ldots\).  Its observation, denoted
\(\run_S(u)\), has one item for every pop in the operation word \(u\):
it records \(\Some{a}\) when type \(a\) is emitted and \(\None\) when the
pop fails.  In the latter case the state is unchanged.
\end{definition}

\begin{definition}[Flush and interleaved equivalence]
For a packet word \(w=a_1\cdots a_n\), put
\[
  \flushOps(w)
  =
  \operatorname{push}(a_1)\cdots\operatorname{push}(a_n)
  \operatorname{pop}^{\,n}.
\]
Schedulers \(S,T\) are \emph{flush-equivalent}, written
\(S\flusheq T\), when
\[
  \run_S(\flushOps(w))=\run_T(\flushOps(w))
  \qquad\text{for every }w\in A^*.
\]
They are \emph{interleaved-equivalent}, written \(S\intereq T\), when
\[
  \run_S(u)=\run_T(u)
  \qquad\text{for every operation word }u,
\]
with pushes and pops arbitrarily interleaved and over-popping observable.
\end{definition}

\begin{theorem}
For every \(k\in\mathbb N\) and every pair of valid stateless schedulers
\(S,T\) over \(\Fin(k)\),
\[
  S\flusheq T \quad\Longrightarrow\quad S\intereq T.
\]
Equivalently, with the definitions expanded, for all \(k,S,T\),
\[
  \begin{gathered}
    \operatorname{Valid}(S)\longrightarrow
    \operatorname{Valid}(T)\longrightarrow\\
    \bigl(\forall w,\ 
      \run_S(\flushOps(w))=\run_T(\flushOps(w))\bigr)
    \longrightarrow\\
    \bigl(\forall u,\ \run_S(u)=\run_T(u)\bigr).
  \end{gathered}
\]
\end{theorem}

\begin{proof}
We establish the equivalent assertion for every finite alphabet \(A\);
renaming its elements along \(A\cong\Fin(|A|)\) then proves the theorem.
Our notation is as follows.  The projection of a word to \(B\subseteq A\)
is \(w|_B\).  For a natural-rank total preorder \(r\),
\(\sort_r(w)\) means the stable rank-sort of \(w\), and \(\pi_P(w)\) is the
word of block names obtained from a partition \(P\).

\medskip
\noindent\textit{Operational invariants.}
Fix a child \(c\) of an internal node.  In any state reachable under valid
paths,
\begin{equation}
 \#\{\text{references to }c\}
 =
 \#\{\text{packets in the subtree below }c\}.
 \tag{1}\label{eq:balance}
\end{equation}
Indeed, the two counts start at zero.  A valid push routed through \(c\)
adds one to each, and a completed pop routed through \(c\) removes one from
each.  Equation \eqref{eq:balance} follows by induction on the run.  In
particular, whenever an internal reference wins, its target contains a
packet.  The recursive failure arm of an atomic pop is therefore never
taken in a reachable state.  A pop returns \(\None\) precisely at total
packet count zero, making the success/failure pattern a function of the
operation word alone.

Only the order of stamps enters a queue comparison.  Thus any
order-preserving relabeling of the stamps leaves the transition sequence
unchanged.  Relabeling a sparse increasing sequence by \(1,2,\ldots\) will
be called \emph{timestamp compression}; it lets us regard a driven child
run with global gaps as a standalone run.

For bookkeeping, give every enqueued entry the identity of the push that
created it.  These ghost identities neither affect a comparison nor alter
the untagged run.  Stamps are injective separately in each queue, and also
in every pending list used later.  Of course one push produces same-stamped
entries at different tree levels; the assertion is only queue-local.
Consequently a selector never has two entries with the same full key
\((\text{rank},\text{arrival})\).

\medskip
\noindent\textit{The two-type calculation.}
Retain only two distinct types \(x,y\) and compare their fixed routes.  Along
the common prefix, every queued value is the same child index.  Whatever
reference is removed there, the node simply asks that child for one packet,
so no type-level choice is made on this prefix.  The first route split, if
there is one, is different: choosing a child chooses between an
\(x\)-only branch and a \(y\)-only branch.  If the routes never split, their
common leaf makes the choice directly.  Collapsing the irrelevant prefix
and the one-type suffixes shows that the restriction is exactly one stable
PIFO with an effective comparison
\[
  e(x,y)\in\{x<y,\ x=y,\ y<x\}.
\]
Writing \(F\) for its flush output, testing the two possible arrival orders
determines which of the three comparisons holds:
\begin{equation}
\begin{array}{c|cc}
e(x,y) & F(xy) & F(yx)\\ \hline
x<y & xy & xy\\
x=y & xy & yx\\
y<x & yx & yx
\end{array}
\tag{2}\label{eq:binary}
\end{equation}
The two columns are both needed: the middle row is the FIFO-equality case.

\medskip
\noindent\textit{Colored-PIFO congruence.}
Consider two selectors on the same pushed occurrences, ordered respectively
by total preorders \(p\) and \(q\), and color each type by \(\gamma\).
Assume exact agreement whenever the colors differ:
\begin{equation}
  \gamma(x)\ne\gamma(y)
  \quad\Longrightarrow\quad
  \cmp_p(x,y)=\cmp_q(x,y).
  \tag{3}\label{eq:cross-agree}
\end{equation}
Equality in \eqref{eq:cross-agree} is part of the assumption.  We claim that
the two selectors produce identical color sequences under any common
sequence of pushes and pops.

Within one color, group types that have the same comparison profile against
all other colors:
\[
 x\sim y
 \quad\Longleftrightarrow\quad
 \gamma(x)=\gamma(y)
 \ \text{ and }\
 \cmp_p(x,z)=\cmp_p(y,z)
 \quad\text{for every }z\text{ of another color}.
\]
Condition \eqref{eq:cross-agree} makes this the same equivalence relation
for \(q\).  Totality supplies the following facts.  First, two distinct
profile classes of one color have a common strict ordering under \(p\) and
\(q\).  Second, if a class ties anything of another color, every member ties
that witness and all other members under both orders; call the class
\emph{pinned}.  Third, for any two classes of different colors, their
comparison is independent of the chosen representatives and is identical
for \(p\) and \(q\).

The coupling records, for each profile class, equal pending cardinalities.
For a pinned class it records more: the two sides have exactly the same
arrival-ordered list of ghost occurrences.  These assertions hold initially
and after a common push.  Consider a pop, and let \(K_p,K_q\) be the classes
containing the two chosen occurrences.  If the classes differ, the count
invariant ensures that a \(K_q\)-occurrence is available to the
\(p\)-selector and a \(K_p\)-occurrence to the \(q\)-selector.  Optimality
then implies
\[
  K_p\le_p K_q,\qquad K_q\le_q K_p,
\]
with either inequality valid for all representatives.  The classes cannot
be two profiles of one color, since their strict ordering is common to both
preorders.  They therefore have different colors, and the constant
cross-class comparison forces a tie.

The tie makes \(K_p\) and \(K_q\) pinned, so their invariant lists put both
chosen occurrences in both queues.  Their rank parts are equal on either
side.  If the choices are \(o_p,o_q\), minimality first gives
\[
 \operatorname{arr}(o_p)\le\operatorname{arr}(o_q)
 \quad\text{and then}\quad
 \operatorname{arr}(o_q)\le\operatorname{arr}(o_p).
\]
The arrivals coincide, and queue-local injectivity says \(o_p=o_q\), which
is impossible for two different classes.  Hence both selectors choose one
profile class and emit one color.  Counts remain matched.  When that class
is pinned, each selector removes the earliest occurrence from the same
list, preserving the stronger clause.  An empty pop is simultaneous because
the total counts match.  The claim follows.

If a use of the lemma is phrased on a subalphabet, it can instead be made
global by padding ranks.  Put every excluded type at one shared rank above
all retained ranks in both orders.  The hypothesis then follows by checking
inside--inside, the two inside--outside orientations, and outside--outside.

\medskip
\noindent\textit{Flush decompositions.}
Suppose \(|A|>1\).  Peel away the initial chain on which every assigned path
names one common child.  Such a node reveals no routing information: each
successful request to it becomes one request to that sole active child,
irrespective of the order among its references.  The chain ends either at
the first genuine branch or at a leaf.  In the leaf case, simulate its queue
by a new selector carrying the leaf's ranks and route each type to its own
one-type FIFO child.  Its root queue mirrors the old leaf queue occurrence
for occurrence, and selecting a reference calls the corresponding
one-type child.

We may therefore replace a scheduler by an online-equivalent object of the
form
\begin{equation}
  M=\bigl(P,r,(M_B)_{B\in P}\bigr),
  \tag{4}\label{eq:normal-form}
\end{equation}
and reuse the name \(M\) for the replacement.  Here \(P\) consists of the
nonempty type sets routed to the active children, \(r\) is the root-rank
preorder, and \(M_B\) is the scheduler below block \(B\).  The partition is
\emph{proper}: it has at least two nonempty blocks, so every \(B\in P\) is
a strict subset of \(A\).

Let \(F_M(w)\) be the emitted type word after pushing all of \(w\) and
draining.  For \(B\in P\),
\begin{equation}
  F_M(w)|_B
  =
  F_{M_B}(w|_B)
  =
  F_M(w|_B).
  \tag{5}\label{eq:subsequence}
\end{equation}
At the start of the drain, child \(B\) has received precisely the projected
input \(w|_B\).  The parent stores one \(B\)-reference for each of those
pushes, so it requests \(\lvert w|_B\rvert\) answers from that child.
Removing the traffic of the other blocks changes neither this local input
nor these requests, proving both equalities.  On the other hand, forgetting
the types and retaining only their root destinations gives
\begin{equation}
  \pi_P(F_M(w))
  =
  \pi_P(\sort_r(w)).
  \tag{6}\label{eq:pattern}
\end{equation}
This is simply the stable ordering of the root entries by \(r\).

We say that \((P,r)\) decomposes \(F\) when these two identities hold.
They imply \emph{block autonomy}.  Namely, if the input uses only \(B\),
the pattern identity keeps the output in \(B\), while
\eqref{eq:subsequence} identifies that output with the child computation.
For unrestricted \(w\), they also give the projection rule
\(F(w)|_B=F(w|_B)\).

\medskip
\noindent\textit{The meet-order lemma.}
The central claim is
\begin{equation}
\boxed{
\begin{gathered}
 (P,r)\text{ and }(Q,s)\text{ decompose }F\\
 \Longrightarrow\ \exists\text{ a total preorder }t\text{ such that}\\
 \cmp_t=\cmp_r\text{ across distinct }P\text{-blocks},\qquad
 \cmp_t=\cmp_s\text{ across distinct }Q\text{-blocks}.
\end{gathered}}
\tag{7}\label{eq:meet-order}
\end{equation}
Let
\[
  R=P\wedge Q
  =
  \{\,B\cap C:\ B\in P,\ C\in Q,\ B\cap C\ne\varnothing\,\}.
\]
For a cell \(D=B\cap C\), the two decomposition identities give
\begin{equation}
\begin{aligned}
 F(w)|_D
 &= (F(w)|_B)|_C                                      \\
 &= F(w|_B)|_C                                         \\
 &= F((w|_B)|_C)                                       \\
 &= F(w|_D).
\end{aligned}
\tag{8}\label{eq:cell-autonomy}
\end{equation}
Thus restriction to any \(D\in R\) commutes with \(F\).

Take \(x,y\) from different cells and define \(e(x,y)\) by the binary table
\eqref{eq:binary}.  If their \(P\)-blocks differ, the \(P\)-normal form
routes their references to different root children.  On inputs using just
\(x,y\), those two destinations determine the emitted type, so the table
reads exactly the root comparison:
\begin{equation}
  \cmp_e(x,y)=\cmp_r(x,y).
  \tag{9}\label{eq:e-r}
\end{equation}
The same reasoning with the \(Q\)-normal form gives, whenever their
\(Q\)-blocks differ,
\begin{equation}
  \cmp_e(x,y)=\cmp_s(x,y).
  \tag{10}\label{eq:e-s}
\end{equation}
When both coordinates differ, the single observable value \(e(x,y)\)
therefore makes the two root comparisons coincide.

We next check triples of meet cells.  Plot a cell \(B\cap C\) at row \(B\)
and column \(C\), and select one type from each of three plotted cells.
Three different rows put every pair under \(r\), and three different
columns put every pair under \(s\), so those cases are already transitive.
If neither occurs, the three occupied squares are the two arms and corner
of an L.  With the corner called \(x\), restricting the two partitions to
the triple yields
\begin{equation}
  \{x,y\}\mid\{z\},
  \qquad
  \{x,z\}\mid\{y\}.
  \tag{11}\label{eq:L}
\end{equation}

Only transitivity can now fail.  If it did, the types could be named so that
\[
  a\le_e b,\qquad b\le_e c,\qquad c<_e a.
\]
Use \(a,b,c\) as both the insertion order and the ensuing flush input.  The
three stamps are different, so arrival order decides every weak root tie.
It suffices to display the orientation in which one partition groups
\(\{b,c\}\) and the other groups \(\{c,a\}\).  On the first root, the two
separated pairs force
\[
  \rho_c<\rho_a\le\rho_b.
\]
Thus the root asks for the blocks in the order
\(\{b,c\},\{a\},\{b,c\}\); if \(\rho_a=\rho_b\), the stamp of \(a\) is
earlier.  The paired child saw \(b\) before \(c\), and \(b\le_e c\), so its
two answers are \(b,c\).  The resulting type word is
\[
  b,a,c.
\]
At the other root the forced inequalities are
\(\sigma_a\le\sigma_b\le\sigma_c\).  Its block requests, in order, are
\[
  \{c,a\},\{b\},\{c,a\}.
\]
The paired child received \(a\) first, but the strict comparison
\(c<_e a\) makes its answers \(c,a\), producing
\[
  c,b,a.
\]
These outputs cannot both equal \(F(abc)\), already because their initial
letters differ.  Rotating \(a,b,c\) moves the grouped edge through its other
two positions; reversing the roles of the partitions covers the three
opposite orientations.  Hence \(e\) restricts to a total preorder on every
choice of three cells.

Encode the cross-cell comparisons in a digraph \(G\).  Insert an arc
\(x\to y\) whenever \(x\le_e y\), and flag it when \(x<_e y\).  Consistency
amounts to showing that no cycle contains a flagged arc.  Put a hypothetical
flagged arc first and call it \(a\to b\); we argue by strong induction on
the cycle length.  If an internal vertex repeats, deleting the intervening
closed segment leaves a shorter flagged cycle, so assume the cycle is
simple.

At length \(2\), the return arc \(b\to a\) is incompatible with
\(a<_e b\).  At length \(3\), every pair of vertices is cross-cell, and the
triple preorder excludes a triangle with a flagged side.

For a length-\(4\) cycle \((a,b,c,d)\), first suppose \(a,c\) are in
distinct cells.  The total preorder on \(a,b,c\), together with
\(a<_e b\le_e c\), makes the chord \(a\to c\) flagged; it closes the
flagged triangle \(a,c,d\).  If instead \(b,d\) are in distinct cells, the total
preorder on \(b,c,d\) supplies \(b\to d\), and \(a,b,d\) is a triangle
whose flagged side is \(a\to b\).  If both diagonals are intra-cell, the
vertices alternate between two cells.  Either \(P\) distinguishes those
cells, in which case all four inequalities are read from \(r\), or it does
not and they are all read from \(s\).  That single rank function cannot
satisfy
\[
  \rho_a<\rho_b\le\rho_c\le\rho_d\le\rho_a.
\]

Now let the length be at least \(5\), with \(c,d\) the next two vertices
after \(b\).  If \(a,c\) are in distinct cells, the three-cell result makes
\(a\to c\) flagged, and replacing \(a\to b\to c\) by this chord gives a
shorter flagged cycle.  If \(b,d\) are in distinct cells, the three-cell
result supplies \(b\to d\); replacing \(b\to c\to d\) preserves the flagged
head \(a\to b\) and again shortens the cycle.  Suppose both diagonals are
intra-cell.  Then \(a,d\) are in distinct cells, since otherwise \(a,c,d\)
would lie in one cell while the consecutive pair \(c,d\) must not.  If the
chord has \(d\to a\) (including a tie), it closes the already excluded
flagged four-cycle \((a,b,c,d)\).  Otherwise the same \(r\)-or-\(s\) rank
function governs all four cross-cell comparisons, and
\[
  a<_e b\le_e c\le_e d;
\]
hence \(a\to d\) is flagged.  Replacing the first three edges by this chord
produces a shorter flagged cycle.  This completes the induction and proves
the claim.

Let \(z\leadsto x\) denote reachability in this graph, including the
length-zero path, and define the natural rank
\begin{equation}
  t(x)
  =
  \#\{\,z:\ z\leadsto x\ \text{and}\ x\not\leadsto z\,\}.
  \tag{12}\label{eq:linearization}
\end{equation}
Along any edge \(x\to y\), the strict-predecessor set counted at \(x\) is
contained in the one counted at \(y\).  If the edge is flagged, \(y\) cannot
reach \(x\), for that would make a flagged cycle; consequently \(x\) is a new
strict predecessor of \(y\), and \(t(x)<t(y)\).  A tied comparison supplies
edges in both directions and gives equal ranks.  Thus the numeric preorder
defined by \(t\) realizes every cross-cell constraint.  By
\eqref{eq:e-r}--\eqref{eq:e-s},
\begin{equation}
  \cmp_t=\cmp_r\text{ across distinct }P\text{-blocks},
  \qquad
  \cmp_t=\cmp_s\text{ across distinct }Q\text{-blocks}.
  \tag{13}\label{eq:t-agreement}
\end{equation}
This proves the meet-order lemma.

\medskip
\noindent\textit{Strong induction on the alphabet.}
Induct simultaneously over the cardinality of a finite alphabet.  With no
types, all pops fail.  With one type, packet count decides whether a pop
succeeds, and there is only one possible successful value.

Now take \(|A|>1\) and flush-equivalent \(S,T\).  Replace both schedulers by
the online-equivalent forms supplied above, keep the names \(S,T\) for those
replacements, and write
\[
  S=\bigl(P,r,(S_B)_{B\in P}\bigr),
  \qquad
  T=\bigl(Q,s,(T_C)_{C\in Q}\bigr).
\]
Their flush transformer is one function \(F\).  Apply the meet-order lemma
to obtain \(R=P\wedge Q\) and \(t\) as in \eqref{eq:t-agreement}.

Change \(S\)'s root ranks from \(r\) to \(t\), leaving all children fixed.
With a block of \(P\) treated as a color, the colored congruence hypothesis
is exactly \eqref{eq:t-agreement}.  The old and new roots consequently call
the same blocks at the same pops, which drives identical operation words
into corresponding children.  If \(S^t\) denotes the changed scheduler,
\[
  S\intereq S^t
  \qquad\text{and}\qquad
  F_{S^t}=F.
\]
The same operation on \(T\) gives
\[
  T\intereq T^t
  \qquad\text{and}\qquad
  F_{T^t}=F.
\]
Both changed schedulers still flush to \(F\).  Their root patterns now use
the same preorder, so \eqref{eq:pattern} says simultaneously that
\begin{equation}
  \pi_P(F(w))=\pi_P(\sort_t(w)),
  \qquad
  \pi_Q(F(w))=\pi_Q(\sort_t(w)).
  \tag{14}\label{eq:both-patterns}
\end{equation}

For each \(D\in R\), fix once and for all a valid scheduler \(U_D\) whose
flush function on \(D\) is \(F|_D\).  Such a scheduler exists: restrict
\(S\) to \(D\) and use \eqref{eq:cell-autonomy}.  Given \(B\in P\), build
\(V_B\) by placing a \(t\)-selector above those \(U_D\) with \(D\subseteq
B\).  We show that this replacement has the same flush function as \(S_B\):
\begin{equation}
  F_{S_B}=F_{V_B}.
  \tag{15}\label{eq:block-flush}
\end{equation}

To identify two \(B\)-words, it is enough to know their word of
\(Q\)-block labels and, separately, their subsequence in every \(Q\)-block.
Indeed, a left-to-right scan of the labels uniquely interleaves those
subsequences.  For \(v\in B^*\), autonomy identifies \(F_{S_B}(v)\) with
\(F(v)\); its \(Q\)-labels are those of \(\sort_t(v)\) by
\eqref{eq:both-patterns}.  The root of \(V_B\) gives exactly those labels as
well, since its children are the nonempty cells \(B\cap C\), \(C\in Q\).

For the remaining subsequences, choose \(C\in Q\) and set \(D=B\cap C\).
When \(D\) is empty, both outputs are \(B\)-valued and have empty
\(C\)-projection.  Otherwise,
\begin{equation}
\begin{aligned}
 F(v)|_C
 &=F(v|_C)\\
 &=F(v|_D)\\
 &=F_{U_D}(v|_D)\\
 &=F_{V_B}(v)|_C.
\end{aligned}
\tag{16}\label{eq:reconstruction}
\end{equation}
The four steps use, in order, \(Q\)-block autonomy, the fact that \(v\) is a
\(B\)-word, the defining choice of \(U_D\), and the child projection formula
for \(V_B\).  Thus label patterns and every labeled subsequence match,
proving \eqref{eq:block-flush}.

Properness of \(P\) gives \(|B|<|A|\), so the induction hypothesis converts
the flush equality into
\begin{equation}
  S_B\intereq V_B.
  \tag{17}\label{eq:child-inter}
\end{equation}
Perform these child replacements simultaneously below the unchanged root
of \(S^t\).  Each \(B\)-child sees the pushes projected to \(B\), and it
receives a pop exactly when the root emits a \(B\)-token.  Equation
\eqref{eq:child-inter} therefore preserves the composite run.  The stamps
in this driven word are sparse global stamps; compressing them supplies the
consecutive stamps assumed by its standalone induction instance.

There are now two consecutive \(t\)-selectors.  Define \(U\) by removing
the middle level and letting one \(t\)-root route directly to the already
fixed children \(U_D\).  Let \(\Omega\) be the pending source pushes, using
their ghost identities.  The coupling invariant can be recorded compactly
as
\[
\begin{aligned}
 \text{state of each two-level }U_D
   &=\text{state of the corresponding one-level }U_D,\\
 \text{occurrences at the outer selector}
   &=\Omega\quad\text{(labeled by \(P\))},\\
 \text{occurrences at \(U\)'s root}
   &=\Omega\quad\text{(labeled by \(R\))},\\
 \text{occurrences at the inner selector for \(B\)}
   &=\Omega|_B\quad\text{(labeled by \(R\))}.
\end{aligned}
\]
Each displayed queue orders occurrences by the identical pair
\((t,\text{arrival})\), whose second component is injective.  Suppose the
outer level chooses \(o\) in block \(B\).  It is the sole minimum of
\(\Omega\); hence it is also the sole minimum of the restricted set
\(\Omega|_B\).  The inner level and \(U\)'s root consequently choose the
same ghost occurrence \(o\), and both invoke the same state \(U_D\).  On
success, all corresponding queues delete \(o\).  On a failure below them,
atomicity restores both composite states.  A push adds the corresponding
entry to each listed queue.  The invariant is preserved, proving the
collapse.  Chaining it with the rank swap and child substitutions gives
\begin{equation}
  S\intereq U.
  \tag{18}\label{eq:S-U}
\end{equation}

Apply the mirror construction to \(T\), crucially retaining the same
chosen family \((U_D)\).  This time the reconstruction takes block labels
from the \(P\)-part of \eqref{eq:both-patterns} and its projected words from
the \(S\)-decomposition.  Each \(C\in Q\) is smaller than \(A\), so
induction and the identical collapse yield
\begin{equation}
  T\intereq U.
  \tag{19}\label{eq:T-U}
\end{equation}
Both original schedulers are therefore interleaved-equivalent to \(U\), and
hence to one another.  Reinsert the failed-pop positions: by
\eqref{eq:balance} these positions depend only on the input operation word.
The resulting equality is equality of the exact lists produced by
\(\run\), including every \(\None\).
\end{proof}

\begingroup
\raggedright
\footnotesize
\paragraph{Kernel-witness scope.}
The colored-PIFO coupling is witnessed by
\texttt{KernelB.lean} (\texttt{Rel}, \texttt{Rel\_pop}) and
\texttt{Kernel.lean} (\texttt{tie\_pins},
\texttt{strict\_\allowbreak of\_\allowbreak sameColor}, \texttt{memInd});
rank padding is witnessed
by \texttt{Refine2A.lean} (\texttt{rankBound}, \texttt{padK},
\texttt{crossAgree\_\allowbreak pad}) and \texttt{Refine2B.lean}
(\texttt{rel\_\allowbreak run\_\allowbreak split}).  The three-cell
calculation is witnessed by
\texttt{Ultra2Flush.lean} (\texttt{flush3\_first},
\texttt{flush3\_mid}, \texttt{flush3\_last}, \texttt{cfg1},
\texttt{cfg2}, \texttt{cfg4},
\texttt{composite\_\allowbreak noCycle}, and
\texttt{composite\_\allowbreak TripleTrans}).  Rank padding is not used in
the main argument above.

The fixed-head marked-cycle construction, including its counting
linearization, is checked by \texttt{KConstruct.lean}
(\texttt{no\_chain}, \texttt{triangle}, \texttt{K}, \texttt{K\_iff}) only
for flat normal forms.  The arbitrary-composite extension is instead
witnessed by
\texttt{Ultra2Extend.\allowbreak preorder\_\allowbreak extension} and
\texttt{Ultra2Flush.\allowbreak composite\_\allowbreak extension}, through
a different
walk-normalization proof.

In \texttt{Ultra2Main.lean}, the root swap is witnessed by \texttt{step} and
\texttt{rankSwap\_\allowbreak interEquiv}.  The reconstruction witnesses
are \texttt{block\_\allowbreak two\_\allowbreak level},
\texttt{block\_\allowbreak flush\_\allowbreak eq}, and
\texttt{reconstruct}; the strict-key collapse witnesses are
\texttt{tmin\_\allowbreak strict} and
\texttt{collapse\_\allowbreak runFrom}, using the pending-list invariant
\texttt{TSorted}.  The compiled capstone derives the two cell families
separately, collapses each side, and then uses cellwise induction; it does
not literally use the preselected common family \((U_D)\) in the prose
organization above.  Finally, \texttt{Stamped.lean} witnesses the
alternative formulation with an explicit arbitrary stamp on every push and
\texttt{sInterEquiv} quantified over arbitrary stamped words; the proof
above instead uses timestamp compression.
\endgroup

\end{document}
